\documentclass[11pt,a4paper]{article}
\usepackage{webclases-lesson}
\lessonSetup{T6}{v0.3}{T6 - Momento lineal, impulso y choques}
\begin{document}
\shorthandoff{<>}
\color{Ink}
\coverHeader{T6 --- Momento lineal, impulso y choques}{Elegir un sistema, estudiar la interacci\'on y comprobar el resultado.}
El momento lineal permite relacionar los estados inmediatamente anteriores y
posteriores a una interacci\'on, incluso si la fuerza cambia mucho durante ella.
La clave es decidir si el impulso externo puede despreciarse.
\begin{lessonindex}
\indexentry{sec:b1}{1. Cambiar el momento: el impulso}
\indexentry{sec:b2}{2. Qu\'e conserva un sistema}
\indexentry{sec:b3}{3. Resolver e interpretar choques}
\end{lessonindex}
\precisionnote

\block{1}{Cambiar el momento: el impulso}{sec:b1}
El momento lineal es $\vec p=m\vec v$: es vectorial, no una energ\'ia.
Sus unidades son $\mathrm{kg\,m/s}$.
Un objeto a $-3\mps$ tiene momento negativo en ese eje,
pero energ\'ia cin\'etica positiva. Para masa constante:
\[
\sum\vec F=\frac{d\vec p}{dt}=m\vec a,\qquad
\vec I=\int_{t_i}^{t_f}\sum\vec F\,dt=\vec p_f-\vec p_i.
\]
El impulso expresa el efecto acumulado de la fuerza durante un intervalo.
$1\units{N\,s}=1\units{kg\,m/s}$.
En una gr\'afica $p(t)$ la pendiente es fuerza; en una gr\'afica $F(t)$ el \'area
es impulso. Los ejes no son intercambiables.

\heading{Fuerza media e instant\'anea}
Para fuerza constante, $\vec I=\vec F\Delta t$. Si cambia durante el contacto,
definimos $\vec F_{\mathrm{med}}=\Delta\vec p/\Delta t$.
Este cociente no informa del valor m\'aximo ni de cada instante.
Una fuerza media peque\~na durante mucho tiempo puede dar el mismo impulso que
una grande durante poco tiempo.

\begin{examplebox}{mismo impulso, distinta fuerza}
\origin{Elaboraci\'on propia.}
\step{Planteamiento} Pelota de $0,200\kg$ que llega con $u=-10,0\mps$
y sale con $v=+5,00\mps$. Calcula el impulso horizontal y la fuerza media
para contactos de $0,0200$ y $0,100\units{s}$.
\solutionblock
\step{Modelo} El eje positivo apunta hacia la salida. Hay inversi\'on de sentido:
$\Delta v$ es $5-(-10)$, no $5-10$.
\step{C\'alculo}
\[
I=m(v-u)=0,200[5,00-(-10,0)]=+3,00\units{N\,s}.
\]
$F_{\mathrm{med}}=3/0,0200=+150\N$ en el contacto breve y
$3/0,100=+30,0\N$ en el prolongado.
\step{Interpretaci\'on} Aumentar el tiempo reduce la fuerza media para el mismo
cambio de momento. No reduce el impulso.
\begin{center}
\begin{tikzpicture}[x=85cm,y=.012cm,font=\small]
\fill[GreenLine!18] (0,0)--(.01,300)--(.02,0)--cycle;
\fill[AmberLine!18] (0,0)--(.05,60)--(.10,0)--cycle;
\draw[axis] (0,0)--(.115,0) node[right] {$t$ (s)};
\draw[axis] (0,0)--(0,345) node[above] {$F_x$ (N)};
\draw[BrandBlue,thick] (0,0)--(.01,300)--(.02,0);
\draw[GreenLine!70!black,dashed,thick] (0,0)--(.05,60)--(.10,0);
\foreach \t in {.02,.05,.10} \draw (\t,0)--(\t,-6) node[below] {$\t$};
\foreach \f in {60,150,300} \draw (0,\f)--(-.001,\f) node[left] {$\f$};
\node[right] at (.025,260) {breve: pico $300\N$};
\node[right] at (.05,100) {prolongado: pico $60,0\N$};
\end{tikzpicture}
\end{center}
Pulsos triangulares \emph{modelizados}, a escala com\'un:
$\tfrac12(0,0200)(300)=\tfrac12(0,100)(60,0)=3,00\units{N\,s}$.
Su fuerza media es la mitad del pico. No son mediciones del contacto.
\variation{Si la pelota se detuviese sin rebotar, $I=0,200[0-(-10)]=+2,00\units{N\,s}$:
ya no ser\'ia el mismo cambio de momento.}
\end{examplebox}

Como comprobaci\'on breve del original, una fuerza \emph{neta} constante de
$5,00\N$ sobre $3,00\kg$ durante $3,00\units{s}$ desde reposo da
$I=15,0\units{N\,s}$ y $v=5,00\mps$ en el sentido de la fuerza.

\block{2}{Qu\'e conserva un sistema}{sec:b2}
El momento total del sistema es $\vec P=\sum_i m_i\vec v_i$.
Las fuerzas internas de una interacci\'on aparecen por pares opuestos sobre
cuerpos distintos. Al sumar las ecuaciones de todos esos cuerpos,
solo las fuerzas externas pueden cambiar el momento total:
\[
\Delta\vec P=\vec I_{\mathrm{ext}}.
\]
Si $\vec I_{\mathrm{ext}}=0$, se conserva $\vec P$. En un choque corto puede
ser una buena aproximaci\'on despreciar el impulso externo horizontal, pero
la brevedad por s\'i sola no lo garantiza. Se compara ese impulso con la escala de
momento relevante. Un cuerpo aislado del resto del choque no conserva su propio
momento: recibe un gran impulso del otro.

\heading{Elegir frontera, direcci\'on e intervalo}
Durante un impacto sobre mesa, peso y normal pueden compensarse verticalmente,
mientras el rozamiento horizontal da un impulso peque\~no. Despu\'es, ese mismo
rozamiento act\'ua durante mucho m\'as tiempo y puede frenar el conjunto.
Por eso se especifican velocidades \emph{inmediatamente} antes y despu\'es del choque.
No se conserva momento entre cualquier par de instantes solo porque haya un choque.

\begin{examplebox}{plastilina y bloque: impacto y frenada}
\origin{Impacto adaptado del original T6, p.\ 3; frenada posterior de elaboraci\'on propia.}
\step{Planteamiento} Plastilina de $0,250\kg$ a $+10,0\mps$ choca con un bloque
de $0,500\kg$ en reposo. Quedan unidos. Desprecia el impulso externo horizontal
durante el impacto. Despu\'es deslizan con $\mu_k=0,20$ y $g=9,81\mpss$.
Halla velocidad com\'un, cambio de $K$ en el choque y distancia hasta parar.
\solutionblock
\step{Modelo} Elegimos ambos cuerpos. La adhesi\'on impone una velocidad final
com\'un; no impone conservaci\'on de $K$.
\step{C\'alculo}
\[
P_i=0,250(10,0)=2,50\units{kg\,m/s},\qquad
v=\frac{2,50}{0,750}=\frac{10}{3}\mps.
\]
$K_i=12,5\J$ y
$K_f=\tfrac12(0,750)(10/3)^2=4,17\J$:
$\Delta K=-8,33\J$.
En la frenada $a=-\mu_kg=-1,96\mpss$ y
$d=v^2/(2\mu_kg)=2,83\units{m}$.
\step{Interpretaci\'on} El momento se conserva aproximadamente en el impacto,
pero pasa de $2,50$ a cero durante la frenada. Su impulso externo posterior es
$-2,50\units{N\,s}$. La energ\'ia cin\'etica que disminuye se transforma
en deformaci\'on, energ\'ia interna y entorno, no desaparece.
\variation{Si el impacto durase $0,0100\units{s}$ y el rozamiento externo fuese
$1,50\N$, su impulso ser\'ia $0,0150\units{N\,s}$, el $0,600\,\text{\%}$
de $P_i$. Esta estimaci\'on ilustra c\'omo justificar la aproximaci\'on;
no es un tiempo de contacto medido.}
\end{examplebox}

\begin{examplebox}{separarse desde reposo}
\origin{Adaptaci\'on del original T6, p.\ 3, con masas y referencia expl\'icitas.}
\step{Planteamiento} Dos patinadores de $60,0$ y $40,0\kg$, juntos y en reposo
sobre hielo ideal, se empujan. El de $60,0\kg$ sale a $+2,00\mps$
respecto al suelo. Halla la otra velocidad y la energ\'ia cin\'etica final.
\solutionblock
\step{Modelo} El sistema son ambos patinadores, sin impulso externo horizontal.
El total inicial es cero: sus momentos finales deben ser opuestos.
\begin{center}
\begin{tikzpicture}[x=1cm,y=1cm,font=\small]
\node at (1,1.4) {antes};
\draw[fill=AmberSoft,draw=AmberLine] (0,0) rectangle (.9,.5);
\draw[fill=BlueSoft,draw=BrandBlue] (.95,0) rectangle (1.85,.5);
\node[below] at (.45,0) {$40\kg$}; \node[below] at (1.4,0) {$60\kg$};
\node at (6.4,1.4) {despu\'es};
\draw[fill=AmberSoft,draw=AmberLine] (4.3,0) rectangle (5.2,.5);
\draw[fill=BlueSoft,draw=BrandBlue] (7,0) rectangle (7.9,.5);
\draw[vectoramber] (4.75,.8)--(3.25,.8) node[midway,above] {$-3,00\mps$};
\draw[vector] (7.45,.8)--(8.45,.8) node[midway,above] {$+2,00\mps$};
\node[below] at (4.75,0) {$40\kg$}; \node[below] at (7.45,0) {$60\kg$};
\end{tikzpicture}
\end{center}
Identidades estables; flechas de velocidad a escala $0,5\units{cm}$ por
$1\mps$, separ\'andose.
\step{C\'alculo}
\[
60(2)+40v_2=0\Rightarrow v_2=-3,00\mps,\qquad
K_f=\tfrac12(60)(2)^2+\tfrac12(40)(3)^2=300\J.
\]
\step{Interpretaci\'on} Los impulsos individuales son $+120$ y
$-120\units{N\,s}$. El momento total sigue siendo cero, pero $K$ aumenta
a costa de energ\'ia interna de los patinadores.
\variation{Si $2,00\mps$ fuese la rapidez de separaci\'on entre ambos,
la condici\'on ser\'ia $v_1-v_2=2,00$, junto a $60v_1+40v_2=0$:
$v_1=+0,800\mps$ y $v_2=-1,20\mps$. La referencia cambia el problema.}
\end{examplebox}

\block{3}{Resolver e interpretar choques}{sec:b3}
En una dimensi\'on fijamos $+x$ a la derecha. $u_1,u_2$ son velocidades
antes del contacto y $v_1,v_2$ despu\'es.
Con impulso externo despreciable:
\[
m_1u_1+m_2u_2=m_1v_1+m_2v_2.
\]
Una sola ecuaci\'on no determina dos velocidades. Necesitamos una condici\'on
del choque: adhesi\'on, conservaci\'on de energ\'ia cin\'etica o restituci\'on.

\heading{Restituci\'on: c\'omo se separan}
Si el cuerpo 1 est\'a a la izquierda y se acerca al 2, $u_1>u_2$.
Definimos el coeficiente de restituci\'on en la direcci\'on del impacto:
\[
e=\frac{v_2-v_1}{u_1-u_2}.
\]
El denominador es la rapidez de acercamiento y el numerador la de separaci\'on.
Para los modelos pasivos tratados:
\begin{tabularx}{\linewidth}{@{}p{2.2cm}Y@{}}
\toprule
Valor & Condici\'on f\'isica\\
\midrule
$e=1$ & Choque el\'astico: se conserva tambi\'en $K$.\\
$0<e<1$ & Inel\'astico: disminuye $K$ y hay separaci\'on.\\
$e=0$ & Igual velocidad final en 1D; perfectamente inel\'astico en esa direcci\'on.\\
\bottomrule
\end{tabularx}
$e=0$ no demuestra por s\'i solo adhesi\'on permanente en todo contacto;
en la plastilina se pegan porque lo dice el enunciado.
La conservaci\'on de energ\'ia \emph{total} no exige conservar solo la parte cin\'etica.

\begin{examplebox}{masas distintas y restituci\'on}
\origin{Adaptado del original T6, p.\ 6; se mantiene $e=0,40$ de forma consistente.}
\step{Planteamiento} $m_1=0,500\kg$, $u_1=+4,00\mps$;
$m_2=0,300\kg$, $u_2=-6,00\mps$, $e=0,40$.
Halla velocidades finales y disminuci\'on de energ\'ia cin\'etica.
\solutionblock
\step{Modelo} Los cuerpos se acercan y el impulso externo es despreciable.
La masa mayor no basta para predecir el sentido final:
hay que comparar momentos con signo.
\[
0,500v_1+0,300v_2=0,200,\qquad
v_2-v_1=0,40[4-(-6)]=4,00.
\]
\step{C\'alculo} Sustituimos $v_2=v_1+4$ en la primera ecuaci\'on:
$0,800v_1+1,20=0,200$.
Por tanto $v_1=-1,25\mps$ y $v_2=+2,75\mps$.
\begin{center}
\begin{tikzpicture}[font=\small]
\foreach \y/\lab in {1.9/antes,0/despu\'es} {
 \node[left] at (-.2,\y) {\lab};
 \draw[fill=BlueSoft,draw=BrandBlue] (1,\y-.25) rectangle (1.7,\y+.25);
 \draw[fill=AmberSoft,draw=AmberLine] (6,\y-.25) rectangle (6.7,\y+.25);
 \node[below] at (1.35,\y-.25) {$m_1$};
 \node[below] at (6.35,\y-.25) {$m_2$};
}
\draw[vector] (1.35,2.4)--++(1.4,0) node[midway,above] {$+4,00$};
\draw[vectoramber] (6.35,2.4)--++(-2.1,0) node[midway,above] {$-6,00$};
\draw[vector] (1.35,.5)--++(-.4375,0) node[left] {$-1,25$};
\draw[vectoramber] (6.35,.5)--++(.9625,0) node[right] {$+2,75$};
\end{tikzpicture}
\end{center}
Velocidades en m/s; escala com\'un $0,35\units{cm}$ por $1\mps$.
La disposici\'on de los cuerpos es esquem\'atica, no una medici\'on de posiciones.
\step{Comprobaci\'on}
$P_f=-0,625+0,825=0,200\units{kg\,m/s}=P_i$.
$K_i=9,40\J$ y $K_f=1,525\J$, que se comunica como $1,53\J$.
La disminuci\'on se calcula antes de redondear:
$9,40-1,525=7,875\J\to7,88\J$.
\step{Interpretaci\'on} Ambos invierten sentido y se separan:
$v_2-v_1=4,00\mps>0$.
\variation{Si quedan unidos, $e=0$ da $v=+0,250\mps$
y $K_f=0,0250\J$. Se conserva el mismo momento con otra condici\'on de choque.}
\end{examplebox}

\heading{Cambiar solo el tipo de choque}
Dos masas iguales de $1,00\kg$ avanzan a $u_1=10,0$ y $u_2=6,00\mps$.
La 1 est\'a detr\'as y alcanza a la 2. En todos los casos,
$P=16,0\units{kg\,m/s}$ y $K_i=68,0\J$.
Resolver $v_1+v_2=16$ y $v_2-v_1=4e$ produce:
\begin{center}
\begin{tabular}{@{}ccccc@{}}
\toprule
$e$ & $v_1$ (m/s) & $v_2$ (m/s) & $K_f$ (J) & Disminuci\'on de $K$ (J)\\
\midrule
0 & 8,00 & 8,00 & 64,0 & 4,00\\
0,50 & 7,00 & 9,00 & 65,0 & 3,00\\
1 & 6,00 & 10,0 & 68,0 & 0\\
\bottomrule
\end{tabular}
\end{center}
Tabla adaptada de la comparaci\'on del original T6 y recalculada.
Las masas iguales intercambian velocidades solo en el choque el\'astico.
Con $m_1=1,00\kg$, $m_2=4,00\kg$ y las mismas velocidades iniciales,
el caso el\'astico da $v_1=3,60$ y $v_2=7,60\mps$: no se intercambian.
Tampoco $e^2$ es una fracci\'on universal de la energ\'ia cin\'etica total conservada:
para $e=0$ en la tabla sigue habiendo $64,0\J$ de energ\'ia cin\'etica.

\begin{exercisebox}{conservar momento no basta}
Dos masas de $1,00\kg$ llegan con $u_1=+4,00\mps$ y $u_2=0$.
Se proponen $v_1=-1,00\mps$ y $v_2=+5,00\mps$ para un choque pasivo.
Calcula $P$, $K$ y $e$; decide si encaja con ese modelo.
\solutionblock
$P_i=P_f=4,00\units{kg\,m/s}$, pero $K$ pasa de $8,00$ a $13,0\J$.
$e=[5-(-1)]/(4-0)=1,50$.
No encaja con un choque pasivo: har\'ia falta liberaci\'on de otra energ\'ia.
$e>1$ no es imposible en cualquier situaci\'on f\'isica; contradice la hip\'otesis
pasiva planteada aqu\'i.
\end{exercisebox}

\heading{Fuentes y procedencia}
Original T6 y background: plastilina, separaci\'on adaptada, restituci\'on y tabla
comparativa. Impulso triangular, frenada posterior y revisi\'on de una soluci\'on
son propuestas propias. Todos los diagramas y c\'alculos se han elaborado de nuevo.
\source{https://fisquiweb.es/Apuntes/apun1B.htm}{FisQuiWeb / IES La Magdalena: inspiraci\'on principal, 1.\textordmasculine\ Bachillerato.}
Contraste externo sobre fuerza media, impulso externo y clasificaci\'on energ\'etica:
\source{https://openstax.org/books/university-physics-volume-1/pages/9-2-impulse-and-collisions}{OpenStax, University Physics 1, 9.2: Impulse and Collisions.}
\source{https://openstax.org/books/university-physics-volume-1/pages/9-3-conservation-of-linear-momentum}{OpenStax, University Physics 1, 9.3: Conservation of Linear Momentum.}
\source{https://openstax.org/books/university-physics-volume-1/pages/9-4-types-of-collisions}{OpenStax, University Physics 1, 9.4: Types of Collisions.}
\end{document}
